% !TEX root = ../Projektdokumentation.tex
\section{Dokumentation}
\label{sec:Dokumentation}
In diesem Abschnitt wird auf die im Rahmen des Projektes erstellte Entwicklerdokumentation eingegangen.


\subsection{Entwicklerdokumentation}
\label{sec:Entwicklerdokumentation}
Da die Anwendung nicht direkt von Nutzern ansprechbar ist benötigt sie keine Anwenderdokumentation.
Lediglich Entwickler müssen wissen, wie man die Anwendung verwendet. Aus diesem Grund wurde eine Entwicklerdokumentation
in Form einer README.md erstellt.
\par\smallskip
Die README.md Datei enthält eine detaillierte Anleitung für die Nutzung der Anwendung bestehend
aus den Voraussetzungen für das erfolgreiche starten, einer Quick Start Anleitung, der Projektstruktur,
einer Auflistung aller verwendeten Konfigurationseinstellungen und Umgebungsvariablen, sowie einer Anleitung für 
das Starten der Applikation mit Docker Compose.
\par\smallskip
Ein Ausschnitt aus der erstellten Benutzerdokumentation befindet sich im \Anhang{app:EntwicklerDoku}. 

Anmerkung: 
Mir fehlt hier ein Hinweis auf weitere Dokumentation im Code, bspw. Summaries an Interfaces, die zu unseren Quellcode Standards gehören. 
Außerdem würde ich vielleicht erwähnen, dass das Projekt in der unternehmensinternen Architekturübersicht aufgenommen wurde (kleine Notlüge, aber nicht ganz falsch).
Weitere Doku zu Projektentscheidungen im Vorfeld befinden sich noch im Confluence, oder? Hier abgrenzen, was aus deiner Feder kommt, aber es ist wichtig, zu erwähnen, 
dass das Projekt nach Abschluss nicht einfach nur existiert.
